\chapter{Board Bring-Up}
\label{chap:fpgalix-board-bring-up}

\section{Introduction}

This chapter brings up the \term{DE1-SoC} board and builds every component the programs of the following chapters depend on. The discussion begins with an overview of the board's boot chain, from reset to a running \gterm{uboot}{U-Boot} (Section~\ref{sec:fpgalix-introduction}), then turns to the more practical side: preparing and building the bootloader itself. Once \gterm{uboot}{U-Boot} is built, the chapter covers the configuration and the building of the \term{Linux} kernel. \gterm{uboot}{U-Boot} loads and boots this kernel, which in turn provides the driver support and core OS services, such as process scheduling, memory management, and the \gterm{v4l2}{V4L2} subsystem later used by the camera driver. The \gterm{buildroot}{Buildroot} root filesystem is built next: a complete embedded \term{Linux} userspace, from the \term{C} library and the \gterm{busybox}{BusyBox}\footnote{\glsentrydesc{busybox}}-provided |init| system and shell utilities, to the user space programs and libraries, \gterm{opencv}{OpenCV}\footnote{\glsentrydesc{opencv}} and \gterm{gstreamer}{GStreamer}\footnote{\glsentrydesc{gstreamer}} among them, needed to build and run \cmd{camera_driver_2} and \cmd{vision_core}. The chapter closes by assembling the bootloader, the kernel, the root filesystem and the \gterm{fpga}{FPGA} bitstream onto a bootable \term{SD} card, producing a working baseline system ready for the examples developed in the remaining chapters.

\begin{table}[htbp]
\centering
\begin{tabularx}{\textwidth}{@{} l X @{}}
\toprule
\textbf{Component} & \textbf{Purpose} \\
\midrule
\gterm{uboot}{U-Boot} & Boots the board: loads the \gterm{fpga}{FPGA} bitstream, brings up the \gterm{hps}{HPS}-to-\gterm{fpga}{FPGA} and \gterm{fpga}{FPGA}-to-\gterm{hps}{HPS} bridges, and starts the \term{Linux} kernel. \\
\term{Linux} kernel & Provides driver support and core OS services, such as process scheduling, memory management and the \gterm{v4l2}{V4L2} subsystem, that the software built in later chapters relies on. \\
\gterm{buildroot}{Buildroot} root filesystem & Provides the embedded \term{Linux} userspace: the \term{init} system, shell utilities, and libraries such as \gterm{opencv}{OpenCV} and \gterm{gstreamer}{GStreamer}, needed to build and run \texttt{camera\_driver\_2} and \texttt{vision\_core}. \\
\gterm{fpga}{FPGA} bitstream & Configures the \gterm{fpga}{FPGA} fabric with the custom hardware described in Chapter~\ref{chap:fpgalix-hw-overview}. \\
\bottomrule
\end{tabularx}
\caption{Main components built in this chapter, and their purpose.}
\label{tab:fpgalix-board-bringup-components}
\end{table}

\section{Required software}

The following components must be installed and available on the \term{PATH}:

\begin{table}[htbp]
\centering
\begin{tabularx}{\textwidth}{@{} l X @{}}
\toprule
\textbf{Software} & \textbf{Why it is needed} \\
\midrule
\term{Ubuntu 22.04 LTS} \term{x86\_64} & Host operating system the build steps in this chapter run on. \\
\gterm{quartusprime}{Quartus Prime} 22.1 & Compiles the \gterm{fpga}{FPGA} hardware project and converts it into the bitstream loaded by \gterm{uboot}{U-Boot} (Section~\ref{sec:fpgalix-bitstream}). \\
\term{ARM GNU} toolchain & Cross-compiles \gterm{uboot}{U-Boot} and the \term{Linux} kernel for the \term{DE1-SoC}'s \term{ARM} \term{Cortex-A9} cores. \\
\bottomrule
\end{tabularx}
\caption{Required host software for this chapter, and why each is needed.}
\label{tab:fpgalix-board-bringup-software}
\end{table}

See Chapter~\ref{chap:fpgalix-devenv} for the installation and configuration of the above, including exact versions, install paths, and the full \gterm{apt}{apt}\footnote{\glsentrydesc{apt}} package list required by the kernel, \gterm{uboot}{U-Boot}, \gterm{buildroot}{Buildroot}, and the programs discussed later in this chapter.

\section{Boot chain overview}
\label{sec:fpgalix-introduction}

On reset, the fixed-function \gterm{bootrom}{Boot ROM} samples the \term{BSEL[2:0]} pins to select the boot source (\term{SD/MMC}, \term{QSPI}, \term{NAND}, or \gterm{fpga}{FPGA} fallback). On the \term{DE1-SoC}, these pins are hardwired on the \gterm{pcb}{PCB} to boot from \term{SD/MMC} \cite{DE1-SCH}, so the \gterm{bootrom}{Boot ROM} reads the combined \gterm{spl}{SPL} (Secondary Program Loader) + \gterm{uboot}{U-Boot} image from a dedicated raw partition (type \term{A2}) on the \term{SD} card, loads it into on-chip \term{RAM}, verifies it, and hands control to \gterm{spl}{SPL}, which brings up clocks, pin muxing and \term{DDR} \cite{CV-TRM}. \gterm{spl}{SPL} then loads full \gterm{uboot}{U-Boot} into \term{DDR}.

From this point, \gterm{uboot}{U-Boot}'s \gterm{bootscript}{boot script} (\repopath{sw/boot.script}, Section~\ref{sec:fpgalix-boot-script}) loads the \gterm{fpgabitstream}{FPGA bitstream} from the \term{SD} card and programs the \gterm{fpga}{FPGA} fabric through the \gterm{hps}{HPS} \gterm{fpgamanager}{FPGA Manager}\footnote{\glsentrydesc{fpgamanager}}, enables the \gterm{hps}{HPS}-to-\gterm{fpga}{FPGA} and \gterm{fpga}{FPGA}-to-\gterm{hps}{HPS} bridges, and then loads the \term{Linux} \gterm{kernel}{kernel} (\texttt{zImage}) together with its \gterm{devicetree}{Device Tree Blob}\footnote{\glsentrydesc{devicetree}}, booting it with \gterm{bootz}{\texttt{bootz}}. The \gterm{kernel}{kernel} then mounts the second \term{SD} card partition (\term{ext4}) as its \gterm{rootfilesystem}{root filesystem} and starts \texttt{init}, which brings up the userspace built by \gterm{buildroot}{Buildroot}.\footnote{See Section~\ref{sec:fpgalix-buildroot}}

\begin{figure}[htbp]
    \centering
    \includesvg[width=\linewidth]{FPGAlixBootChain}
    \caption{\term{HPS} boot chain, from reset to the \term{Linux} kernel.}
    \label{fig:fpgalix-boot-chain}
\end{figure}

\section{Bootloader generation flow}
\label{sec:fpgalix-bootloader-flow}

\term{SoC EDS} (Embedded Design Suite) was \gterm{altera}{Altera}'s, later \term{Intel}'s, companion toolset to \term{Quartus} for \term{SoC} \gterm{fpga}{FPGA} embedded development. Since \term{SoC EDS} v19.1 Std / v19.3 Pro, the bootloader (\gterm{uboot}{U-Boot}) build flow no longer relies on its \gterm{gui}{GUI}\footnote{\glsentrydesc{gui}} tool, \cmdpath{bsp-editor}, to configure and generate \gterm{uboot}{U-Boot}. The current flow proceeds as follows:

\begin{enumerate}
    \item The hardware project is compiled in \gterm{quartusprime}{Quartus Prime}: this produces a handoff folder, a set of \term{XML} files describing how the \gterm{hps}{HPS} is configured (pin muxing, clocks, \term{SDRAM} timings, etc.).
    \item \cmd{cv_bsp_generator.py} \cite{MACNICA-BOOT} (part of the \gterm{uboot}{U-Boot} source tree) is run on the handoff folder, generating the board-specific header files (\repopath{pinmux_config.h}, \repopath{pll_config.h}, \repopath{sdram_config.h}, \repopath{iocsr_config.h}).
    \item These headers are copied into \gterm{uboot}{U-Boot}'s generic source code, together with any custom user options (device tree, defconfig, etc.).
    \item \gterm{make}{\cmdpath{make}} is invoked to build the bootloader image.
\end{enumerate}

\begin{figure}[htbp]
    \centering
    \includesvg[width=\linewidth]{FPGAlixBootGenerationFlow}
    \caption{Current bootloader generation flow (\term{SoC EDS} v19.1 Std / v19.3 Pro and later). Gray boxes are \term{Quartus Prime}; green boxes are part of \term{U-Boot}.}
    \label{fig:fpgalix-boot-generation-flow}
\end{figure}

\begin{warningbox}
    \gterm{platformdesigner}{Platform Designer}'s \term{IP} search path (Tools > Options > IP Catalog, ``IP Search Path'') is a \term{Quartus}-wide setting, not scoped to a single project. Every custom component used by this design (\hw{pclk_reset_controller}, \hw{ov7670_data_interface}, and the remainder found under \repopath{hw/my_vhdl/qsys_components/}, Chapter~\ref{chap:fpgalix-hw-overview}) resides there, and \gterm{platformdesigner}{Platform Designer} is able to resolve them only if that folder is listed. This setting should therefore be verified not only when configuring a new machine for the first time, but on every occasion where development switches between \prjname{FPGAlix} and another \gterm{platformdesigner}{Platform Designer} project on the same machine: each project maintains its own \repopath{hw/my_vhdl/qsys_components/}, and a search path left pointing at the other project's folder causes step 1 above to fail.
\end{warningbox}

\section{Getting the repository}

The remainder of this chapter assumes that the \prjname{FPGAlix} repository has already been cloned locally, as follows, including submodules:

\begin{shellcodehost}
git clone --recurse-submodules https://github.com/tceccarini/FPGAlix.git
cd FPGAlix
\end{shellcodehost}

The |--recurse-submodules| flag causes \gterm{git}{Git}\footnote{\glsentrydesc{git}} to additionally clone every submodule belonging to the project at the same time; these are located under \repopath{repos/} (including |u-boot-socfpga| \cite{SUB-UBOOT}, required in the following section, and |linux-socfpga| \cite{SUB-LINUX} and |buildroot| \cite{SUB-BUILDROOT}, required in Sections~\ref{sec:fpgalix-kernel} and \ref{sec:fpgalix-buildroot}). In its absence, the corresponding directories would be created but left empty.

\section{Building the bootloader}
\label{sec:fpgalix-bootloader}

The following commands implement the flow shown in Figure~\ref{fig:fpgalix-boot-generation-flow}, and are to be executed from the root of the repository.

\subsection{Generate the BSP headers from the Quartus handoff}

Compilation of the hardware project causes \term{Quartus} to export the \gterm{hps}{HPS} configuration as \term{XML} files under the handoff folder \repopath{hw/quartus/hps_isw_handoff/soc_system_hps_0/}. This configuration is then converted into the \term{C} header files required by \gterm{uboot}{U-Boot} by means of the following command:

\begin{shellcodehost}
python3 repos/u-boot-socfpga/arch/arm/mach-socfpga/cv_bsp_generator/cv_bsp_generator.py
    -i hw/quartus/hps_isw_handoff/soc_system_hps_0/
    -o repos/u-boot-socfpga/board/terasic/de1-soc/qts/
\end{shellcodehost}

This step must be repeated whenever the \gterm{hps}{HPS} subsystem is modified in \gterm{platformdesigner}{Platform Designer}; otherwise, \gterm{uboot}{U-Boot} initializes the hardware according to stale configuration data.

\subsection{Configure and build U-Boot}

\begin{shellcodehost}
cd repos/u-boot-socfpga
export ARCH=arm
export CROSS_COMPILE=arm-none-linux-gnueabihf-
make socfpga_cyclone5_defconfig
make menuconfig
\end{shellcodehost}

Within |menuconfig|, select the board target as follows: \term{ARM} architecture $\rightarrow$ \gterm{altera}{Altera} \term{SOCFPGA} board select $\rightarrow$ \term{Terasic} \term{DE1-SoC} (\term{Cyclone V}).

Upon exiting, confirm \term{Yes} when prompted to save the configuration; failure to do so leaves the incorrect board target selected, and the build silently targets the wrong hardware.

\begin{shellcodehost}
make -j$(nproc)
\end{shellcodehost}

\subsection{Result}

The build produces \repopath{repos/u-boot-socfpga/u-boot-with-spl.sfp}, an image combining \gterm{spl}{SPL} and \gterm{uboot}{U-Boot} in the \gterm{altera}{Altera} \term{SoCFPGA} format. This file is subsequently written, in raw form, to the boot partition of the \term{SD} card; it is also accessible via the symlink \repopath{sdcard/u-boot-with-spl.sfp}.

\section{Building the Linux kernel}
\label{sec:fpgalix-kernel}

Enter the \term{Linux} source directory, and set the required environment variables:

\begin{shellcodehost}
cd repos/linux-socfpga
export ARCH=arm
export CROSS_COMPILE=arm-none-linux-gnueabihf-
make socfpga_defconfig
\end{shellcodehost}

\subsection{Kernel configuration}

Open the interactive configuration menu with:

\begin{shellcodehost}
make menuconfig
\end{shellcodehost}

\subsubsection{Base configuration}

The base kernel configuration required by this project is as follows:

\begin{itemize}
    \item \texttt{Device Drivers} $\rightarrow$ \texttt{DMA Engine support}
    \begin{itemize}
        \item \texttt{Altera / Intel mSGDMA Engine} $\rightarrow$ * (built-in)
    \end{itemize}
    \item \texttt{Device Drivers} $\rightarrow$ \texttt{Input device support} $\rightarrow$ \texttt{Keyboards}
    \begin{itemize}
        \item \texttt{GPIO Buttons} $\rightarrow$ * (built-in)
    \end{itemize}
    \item \texttt{Device Drivers} $\rightarrow$ \texttt{LED Support}
    \begin{itemize}
        \item \texttt{LED Class Support} $\rightarrow$ * (built-in)
        \item \texttt{LED Trigger support} $\rightarrow$ \texttt{LED Heartbeat Trigger} $\rightarrow$ * (built-in)
    \end{itemize}
\end{itemize}
\subsubsection{Configuration required by camera\_driver\_2}
The following options are required specifically by \repopath{sw/camera_driver_2}:\footnote{Its own \texttt{README} documents the rationale for each; they are consolidated here so that the kernel need only be built once.}

\begin{itemize}
    \item \texttt{Device Drivers} $\rightarrow$ \texttt{Multimedia support}
    \begin{itemize}
        \item \texttt{Multimedia support} $\rightarrow$ * (built-in, |CONFIG_MEDIA_SUPPORT|). The option ``Filter media drivers'' should be left unchecked, as a consequence of which |MEDIA_CAMERA_SUPPORT| and |VIDEO_DEV| are enabled automatically.
    \end{itemize}
    \item \texttt{Device Drivers} $\rightarrow$ \texttt{Multimedia support} $\rightarrow$ \texttt{Media drivers} $\rightarrow$ \texttt{V4L platform devices}
    \begin{itemize}
        \item \texttt{Aspeed AST2400/AST2500 Video Engine} $\rightarrow$ * (built-in, |CONFIG_VIDEO_ASPEED|). |CONFIG_VIDEOBUF2_DMA_CONTIG| does not possess a corresponding |menuconfig| entry, being a hidden symbol; enabling this otherwise unrelated platform driver causes it to be selected as a side effect, since the driver itself depends on it. It is \extpath{videobuf2-dma-contig.ko} that is actually required by \cmd{camera_driver_2}; the \term{Aspeed} driver itself is never loaded on the target.
    \end{itemize}
    \item \texttt{Device Drivers} $\rightarrow$ \texttt{I\textsuperscript{2}C support} $\rightarrow$ \texttt{I\textsuperscript{2}C Hardware Bus support}
    \begin{itemize}
        \item \texttt{Altera Soft IP I\textsuperscript{2}C} $\rightarrow$ * (built-in, |CONFIG_I2C_ALTERA|). In the absence of this option, the \hw{altera_avalon_i2c} \gterm{platformdesigner}{Platform Designer} component (\texttt{compatible = "altr,softip-i2c-v1.0"}) remains invisible to \term{Linux}, preventing the camera from being configured over \gterm{i2c}{I\textsuperscript{2}C} whenever the bus is routed through the \gterm{fpga}{FPGA} rather than the \gterm{hps}{HPS}.
    \end{itemize}
    \item \texttt{Memory Management options}
    \begin{itemize}
        \item \texttt{Contiguous Memory Allocator} $\rightarrow$ * (|CONFIG_CMA|)
    \end{itemize}
    \item \texttt{Library routines}
    \begin{itemize}
        \item \texttt{DMA Contiguous Memory Allocator} $\rightarrow$ * (|CONFIG_DMA_CMA|), \texttt{Size in Mega Bytes} $\rightarrow$ 16. This setting is recommended rather than strictly mandatory: the requirement of three buffers of 640$\times$480$\times$1 byte amounts to under 1 MB, so the 16 MB default provides an ample margin.
    \end{itemize}
\end{itemize}

Upon exiting |menuconfig|, confirm \term{Yes} when prompted to save.

Build the compressed kernel image and prepare the tree for external kernel modules:

\begin{shellcodehost}
make -j$(nproc) zImage
make modules_prepare
make -j$(nproc) modules
\end{shellcodehost}

\subsection{Device Tree Source}

Prior to building the \gterm{devicetree}{Device Tree} blobs, a board-specific \gterm{devicetree}{DTS} file must exist at:

\begin{shellcodehost}
repos/linux-socfpga/arch/arm/boot/dts/socfpga_cyclone5_fpgalix.dts
\end{shellcodehost}

This file describes the hardware topology of the \term{SoC} and is maintained manually as \repopath{socfpga_cyclone5_fpgalix.dts}, which includes \repopath{socfpga_cyclone5.dtsi} (the standard \term{SoCFPGA} base provided by the kernel) and adds the board-specific nodes on top of it.

\subsubsection{The Device Tree Source is already provided}

\repopath{socfpga_cyclone5_fpgalix.dts} is not part of the upstream kernel; it was authored for this project and is already committed to this project's \term{Linux} kernel fork (|linux-socfpga|), which arrives as a submodule, matching the current hardware. If the \gterm{platformdesigner}{Platform Designer} hardware configuration has not changed, this file requires no further editing.

\subsubsection{Compiling the Device Tree Blob}

Unlike the \gterm{devicetree}{DTS} source, the compiled \gterm{devicetree}{Device Tree Blob} (\repopath{.dtb}) is a build artifact: it is not tracked by the submodule, it does not exist after a fresh clone, and it must be compiled locally on every build, regardless of whether the \gterm{devicetree}{DTS} source itself was changed or not. This step is therefore always required, even when Paragraph~\ref{subsec:fpgalix-dts-update} is skipped entirely:

\begin{shellcodehost}
make dtbs
\end{shellcodehost}

The compiled \gterm{devicetree}{DTB} is accessible via the symlink \repopath{sdcard/boot_partition/socfpga.dtb}, which points to \repopath{repos/linux-socfpga/arch/arm/boot/dts/socfpga_cyclone5_fpgalix.dtb}.

\subsection{Updating the Device Tree Source}
\label{subsec:fpgalix-dts-update}

This section only applies when the \gterm{hps}{HPS} subsystem has been modified in \gterm{platformdesigner}{Platform Designer}; otherwise \repopath{socfpga_cyclone5_fpgalix.dts} already matches the current hardware, and this section can be skipped.

\subsubsection{Using sopc2dts as a reference}

|sopc2dts| is \gterm{altera}{Altera}'s own tool for generating a device tree from a |.sopcinfo| file, tracked in this project as the \repopath{repos/sopc2dts} submodule (a fork of \extpath{altera-opensource/sopc2dts}). It cannot, however, be employed as a direct \gterm{devicetree}{DTS} source for the kernel: its output is a standalone file that does not include \repopath{socfpga_cyclone5.dtsi} and is therefore incompatible with the modern \term{Linux} \term{SoCFPGA} \gterm{devicetree}{DTS} infrastructure.

It remains useful, however, as a reference: its output reflects the current \gterm{platformdesigner}{Platform Designer} configuration, and is compared against \repopath{socfpga_cyclone5_fpgalix.dts} to identify which addresses, interrupt numbers and compatible strings changed, updating the latter accordingly.

The \repopath{sopc2dts} utility itself is not prebuilt: it must be compiled from source once, the first time it is used:

\begin{shellcodehost}
cd repos/sopc2dts
make all
\end{shellcodehost}

Generate the reference output with:

\begin{shellcodehost}
java -jar sopc2dts.jar --input ../../hw/quartus/soc_system.sopcinfo --output /tmp/socfpga_ref.dts
\end{shellcodehost}

\subsubsection{Authoring the Device Tree Source from the reference}

In this project, \repopath{socfpga_cyclone5_fpgalix.dts} was authored from the \repopath{sopc2dts} reference output using \gterm{claudecode}{Claude Code}, an \term{AI} coding assistant, rather than written by hand. The reference \gterm{devicetree}{DTS} generated by \repopath{sopc2dts} (\extpath{/tmp/socfpga_ref.dts} above) was given as input, not the raw \repopath{soc_system.sopcinfo}, generated by \term{Quartus}: providing the latter directly led to significantly less accurate results, since it left the extraction of addresses, interrupt numbers and compatible strings to the \term{AI} assistant itself, rather than to \repopath{sopc2dts}. Manual authoring remains entirely possible: \repopath{socfpga_cyclone5_fpgalix.dts} is a plain text file, and the task consists solely of transcribing addresses, interrupt numbers and compatible strings from the reference output, for those who prefer to do so directly.

\section{Building the root filesystem (Buildroot)}
\label{sec:fpgalix-buildroot}

Enter the \gterm{buildroot}{Buildroot} directory and set the required environment variables:

\begin{shellcodehost}
cd repos/buildroot
export ARCH=arm
export CROSS_COMPILE=arm-none-linux-gnueabihf-
\end{shellcodehost}

\subsection{Buildroot configuration}
\label{subsec:fpgalix-buildroot-config}

Open the configuration menu with:

\begin{shellcodehost}
make menuconfig
\end{shellcodehost}

\subsubsection{Base configuration}
\label{subsec:fpgalix-buildroot-base}

The base configuration required by this project is as follows:

\begin{itemize}
    \item \texttt{Target options}
    \begin{itemize}
        \item \texttt{Target Architecture} $\rightarrow$ \term{ARM} (little endian)
        \item \texttt{Target Architecture Variant} $\rightarrow$ \term{cortex-A9}
        \item \texttt{Enable NEON SIMD extension support} $\rightarrow$ Y
        \item \texttt{Enable VFP extension support} $\rightarrow$ Y
        \item \texttt{Floating point strategy} $\rightarrow$ \gterm{neon}{NEON}\footnote{\glsentrydesc{neon}}
    \end{itemize}
    \item \texttt{Toolchain}
    \begin{itemize}
        \item \texttt{C library} $\rightarrow$ \term{glibc}
        \item \texttt{Enable C++ support} $\rightarrow$ Y
        \item \texttt{Build cross gdb for the host} $\rightarrow$ Y
        \item \texttt{Target ABI} $\rightarrow$ \term{EABIhf}
    \end{itemize}
    \item \texttt{System configuration}
    \begin{itemize}
        \item \texttt{/dev management} $\rightarrow$ Dynamic using \term{devtmpfs} + \term{eudev}
        \item \texttt{Enable root login with password} $\rightarrow$ Y
        \item \texttt{Root password} $\rightarrow$ (set a password)
    \end{itemize}
    \item \texttt{Target packages} $\rightarrow$ \texttt{Debugging, profiling and benchmark}
    \begin{itemize}
        \item \gterm{gdb}{gdb}\footnote{\glsentrydesc{gdb}} $\rightarrow$ Y, then enable \gterm{gdbserver}{gdbserver}\footnote{\glsentrydesc{gdbserver}}
        \item \term{strace} $\rightarrow$ Y
    \end{itemize}
    \item \texttt{Target packages} $\rightarrow$ \texttt{Hardware handling}
    \begin{itemize}
        \item \texttt{Show packages that are also provided by busybox} $\rightarrow$ Y
        \item \term{i2c-tools} $\rightarrow$ Y
    \end{itemize}
    \item \texttt{Target packages} $\rightarrow$ \texttt{Networking applications}
    \begin{itemize}
        \item \term{chrony} $\rightarrow$ Y (NTP client)
        \item \gterm{dhcpcd}{dhcpcd} $\rightarrow$ Y (DHCP client)
        \item \term{ethtool} $\rightarrow$ Y (Ethernet diagnostics)
        \item \term{ifupdown-scripts} $\rightarrow$ Y (ifup/ifdown support)
        \item \term{iperf3} $\rightarrow$ Y (bandwidth testing)
        \item \term{iproute2} $\rightarrow$ Y (ip addr, ip route)
        \item \gterm{openssh}{openssh} $\rightarrow$ Y (SSH server and client)
        \item \term{sshfs} $\rightarrow$ Y (mount remote filesystems over SSH)
        \item \term{tcpdump} $\rightarrow$ Y (network packet capture)
    \end{itemize}
    \item \texttt{Target packages} $\rightarrow$ \texttt{System tools}
    \begin{itemize}
        \item \term{htop} $\rightarrow$ Y
        \item \term{kmod} $\rightarrow$ Y
        \item \term{util-linux} $\rightarrow$ Y
        \item \term{xz} $\rightarrow$ Y
        \item \term{zip} $\rightarrow$ Y
    \end{itemize}
    \item \texttt{Target packages} $\rightarrow$ \texttt{Text editors and viewers}
    \begin{itemize}
        \item \term{nano} $\rightarrow$ Y
    \end{itemize}
\end{itemize}

This last setting is not selected automatically: enabling ``Enable VFP extension support'' makes several floating-point strategies available (\term{VFPv3}, \gterm{neon}{NEON}, among others), and since the default chosen by \gterm{buildroot}{Buildroot} is not \gterm{neon}{NEON}, it must be selected explicitly. The \gterm{neon}{NEON}-accelerated \gterm{opencv}{OpenCV} and \term{libjpeg-turbo} code of \repopath{sw/vision_core} depends on this setting.

\subsubsection{Configuration required by camera\_driver\_2 and vision\_core}

The following packages, in addition to those already listed, are required specifically by \repopath{sw/camera_driver_2} and \repopath{sw/vision_core}:\footnote{Their respective \texttt{README}s document the rationale for each; they are consolidated here so that a single build suffices.}

\begin{itemize}
    \item \texttt{Target packages} $\rightarrow$ \texttt{Libraries} $\rightarrow$ \texttt{Hardware handling}
    \begin{itemize}
        \item \term{libv4l} $\rightarrow$ Y, with \term{v4l-utils} tools enabled (provides \term{v4l2-ctl})
    \end{itemize}
    \item \texttt{Target packages} $\rightarrow$ \texttt{Libraries} $\rightarrow$ \texttt{Graphics}
    \begin{itemize}
        \item \term{opencv4} $\rightarrow$ Y, with the following enabled: \term{imgcodecs}, \term{imgproc}, \term{videoio}, \term{jpeg} support, \term{v4l} support, \term{tbb} support
        \item \term{libjpeg-turbo} $\rightarrow$ Y
    \end{itemize}
    \item \texttt{Target packages} $\rightarrow$ \texttt{Audio and video applications}
    \begin{itemize}
        \item \gterm{gstreamer}{gstreamer} 1.x $\rightarrow$ Y, with the following sub-plugins enabled:
        \begin{itemize}
            \item \term{gst1-plugins-base} $\rightarrow$ app
            \item \term{gst1-plugins-base} $\rightarrow$ videoconvertscale
            \item \term{gst1-plugins-good} $\rightarrow$ rtpmanager (not rtp)
            \item \term{gst1-plugins-bad} $\rightarrow$ openh264 (\gterm{h264}{H.264}\footnote{\glsentrydesc{h264}} output only)
            \item \term{gst1-rtsp-server} $\rightarrow$ Y
        \end{itemize}
    \end{itemize}
\end{itemize}

Upon exiting |menuconfig|, confirm \term{Yes} when prompted to save.

\subsection{Build}
\label{subsec:fpgalix-buildroot-build}

\begin{shellcodehost}
make
\end{shellcodehost}

Do not use the |-j$(nproc)| flag: \gterm{buildroot}{Buildroot} manages parallelism internally. The default settings already instruct \gterm{buildroot}{Buildroot} to employ all available \term{CPU} cores throughout the entire build (toolchain, packages, and root filesystem).

The root filesystem image is generated at \repopath{repos/buildroot/output/images/rootfs.tar}, accessible via the symlink \repopath{sdcard/rootfs.tar}, alongside the remaining \term{SD} card artifacts.

This same build also produces the toolchain used by the \gterm{make}{\cmdpath{make arm}} target of \repopath{sw/vision_core}; this \gterm{buildroot}{Buildroot} build must therefore complete before the \cmd{vision_core} program can be cross-compiled.

\section{Generating the FPGA bitstream (RBF)}
\label{sec:fpgalix-bitstream}

The \gterm{fpga}{FPGA} fabric on the \term{DE1-SoC} may be configured through several paths: over \gterm{jtag}{JTAG} from \term{Quartus}, over \gterm{jtag}{JTAG} via a \gterm{jtag}{JTAG} adapter tool, or by the \gterm{hps}{HPS} itself at boot, through the \gterm{fpgamanager}{FPGA Manager}. \prjname{FPGAlix} employs the last of these: \gterm{uboot}{U-Boot} loads the bitstream from the \term{SD} card and programs the fabric through the \gterm{fpgamanager}{FPGA Manager}, prior to booting \term{Linux}.

\subsection{MSEL configuration (SW10)}
\label{subsec:fpgalix-msel}

Which of these paths is used is a hardware setting, selected by the \term{MSEL[4:0]} pins. For \gterm{uboot}{U-Boot} to be the one that programs the \gterm{fpga}{FPGA} fabric, as described above, \term{MSEL[4:0]} must select the \gterm{fpgamanager}{FPGA Manager} path. On the \term{DE1-SoC}, this is set with the \term{SW10} \term{DIP} switch, located on the underside of the board.

\textbf{Note:} \term{SW10} employs negative logic: the \term{ON} position corresponds to logic 0, and the \term{OFF} position to logic 1.

Set the six \term{SW10} switches, on the underside of the board, to the positions given in the table below, which select the \gterm{fpgamanager}{FPGA Manager} path (\term{MSEL[4:0]} = 00000):

\begin{table}[htbp]
\centering
\begin{tabular}{@{} l l l l @{}}
\toprule
\textbf{Switch} & \textbf{Signal} & \textbf{Position} & \textbf{Logic} \\
\midrule
\term{SW10.1} & \term{MSEL0} & \term{ON} & 0 \\
\term{SW10.2} & \term{MSEL1} & \term{ON} & 0 \\
\term{SW10.3} & \term{MSEL2} & \term{ON} & 0 \\
\term{SW10.4} & \term{MSEL3} & \term{ON} & 0 \\
\term{SW10.5} & \term{MSEL4} & \term{ON} & 0 \\
\term{SW10.6} & --- & \term{OFF} & 1 \\
\bottomrule
\end{tabular}
\caption{\term{SW10} switch positions for the \term{FPGA Manager} path (\term{MSEL[4:0]} = 00000).}
\label{tab:fpgalix-sw10-msel}
\end{table}

\subsection{Converting the SOF to RBF}

\term{Quartus} compiles the hardware project into a \gterm{sof}{\repopath{.sof}} (\term{SRAM Object File}), the complete \gterm{fpga}{FPGA} configuration, expressed in a format understood only by \term{Quartus}/\gterm{jtag}{JTAG} tools. The \gterm{fpgamanager}{FPGA Manager} mechanism, used by \gterm{uboot}{U-Boot}, instead requires an \gterm{rbf}{\repopath{.rbf}} (\term{Raw Binary File}), a stripped-down version of the same configuration, obtained by conversion from the root of the repository:

\begin{shellcodehost}
quartus_cpf -c hw/quartus/output_files/FPGAlix.sof sdcard/boot_partition/soc_system.rbf
\end{shellcodehost}

The resulting output is placed directly under \repopath{sdcard/boot_partition/}, alongside the remaining files destined for the \term{SD} card's boot partition.

\section{Boot script}
\label{sec:fpgalix-boot-script}

\repopath{sw/boot.script} was authored for this project: a human-readable \gterm{uboot}{U-Boot} script that defines the boot sequence. It persists the automatically generated \term{Ethernet MAC} address on first boot, sets the kernel command line (serial console at 115200 baud, root filesystem on the \term{SD} card's second partition, mounted read-write), loads and programs the \gterm{fpga}{FPGA} bitstream generated in Section~\ref{sec:fpgalix-bitstream}, enables the \gterm{hps}{HPS}-to-\gterm{fpga}{FPGA} and \gterm{fpga}{FPGA}-to-\gterm{hps}{HPS} bridges, and finally loads the \term{Linux} kernel (\texttt{zImage}) and the \gterm{devicetree}{Device Tree Blob} (\repopath{socfpga.dtb}), booting the system.

Compile this script into the binary image that \gterm{uboot}{U-Boot} expects to find on the boot partition:

\begin{shellcodehost}
cd sw/
mkimage -A arm -O linux -T script -C none -a 0 -e 0 -n "Boot" -d boot.script ../sdcard/boot_partition/u-boot.scr
\end{shellcodehost}

This step must be repeated whenever \repopath{sw/boot.script} is modified.

\section{Preparing the SD card}

\begin{warningbox}
    This section involves low-level disk operations. An incorrect device name or a mistyped command can silently overwrite a system disk, a data partition, or any other block device on the machine in use. Every device path (\extpath{/dev/sdX}, \extpath{/dev/loopN}) should be double-checked before pressing Enter; when in doubt, the operation should be stopped and verified.
\end{warningbox}

The \term{SD} card requires three partitions: a \term{FAT32} boot partition (kernel, \gterm{devicetree}{DTB}, bitstream, boot script), an \term{ext4} root filesystem, and the raw \term{A2} partition for \gterm{uboot}{U-Boot}. All files destined for the card are collected under \repopath{sdcard/}. A \gterm{diskimagefile}{disk image file} is constructed first and subsequently written to the physical \term{SD} card; the same procedure may be applied directly to the \term{SD} card by substituting the \term{SD} card device (e.g. \extpath{/dev/sdX}) for \repopath{sdcard/images/sdcard.img}, thereby omitting the image creation step and the final \gterm{dd}{\cmdpath{dd}}/\term{Etcher} write.

The final image may be written to the \term{SD} card using:

\begin{itemize}
    \item \term{Linux}: \gterm{dd}{\cmdpath{dd}} (see the last step below)
    \item \term{Windows}/\term{Mac}: \gterm{balenaetcher}{balenaEtcher}\footnote{\gterm{balenaetcher}{balenaEtcher} --- \url{https://etcher.balena.io/}}, free and open-source software; the image and target drive are selected accordingly.
\end{itemize}

\subsection{Create the image file}

From the root of the repository, enter the \repopath{sdcard/} directory:

\begin{shellcodehost}
cd sdcard/
mkdir -p images
\end{shellcodehost}

Create an empty image, sized to accommodate any ``8 GB'' \term{SD} card. \term{SD} card manufacturers employ decimal units (1 GB = 1{,}000{,}000{,}000 bytes), such that an ``8 GB'' card provides approximately 7629 \term{MiB} of actual space; 7400 \term{MiB} is used here to leave a safe margin:

\begin{shellcodehost}
dd if=/dev/zero of=images/sdcard.img bs=1M count=7400
\end{shellcodehost}

\subsection{Partition the image}

\begin{shellcodehost}
fdisk images/sdcard.img
\end{shellcodehost}

Within \cmdpath{fdisk}, enter the following commands in order:

\begin{shellcodehost}
o       # new MBR partition table
n       # new partition 1, FAT32 boot (500 MiB)
p       
1
(enter)
+500M
t       # change partition type
c       # set type to W95 FAT32 LBA (partition 1, FAT32 boot, 500 MiB)
p       # print disk info, note "Sectors" and "Sector size"
\end{shellcodehost}

Prior to creating partition 3, calculate its first sector:

\[
\text{first\_sector}_{p3} = \text{total\_sectors} - \left\lceil \frac{10 \times 1024 \times 1024}{\text{sector\_size}} \right\rceil
\]

Example, assuming 512-byte sectors and a 7400 \term{MiB} image (total sectors = 15{,}155{,}200):

\[
\begin{aligned}
\text{first\_sector}_{p3} &= 15{,}155{,}200 - \left\lceil \frac{10 \times 1024 \times 1024}{512} \right\rceil \\
&= 15{,}155{,}200 - 20{,}480 \\
&= 15{,}134{,}720
\end{aligned}
\]

Create partition 3 at that position:

\begin{shellcodehost}
n               # partition 3, raw A2, U-Boot SPL (10 MiB)
p
3
15134720        # enter the calculated first sector
(enter)         # accept default last sector (end of disk), fdisk should confirm ~10 MiB
t
3
a2              # set type to A2 (Altera SoCFPGA bootloader)
n
p
2
(enter)         # fdisk auto-selects first sector after p1
(enter)         # fdisk auto-selects last sector before p3
p               # verify the partition table before writing
w               # write and exit
\end{shellcodehost}

For this example (512-byte sectors, 7400 \term{MiB} image), the expected output of \ccode{p} is as follows:

\begin{shellcodehost}
Disk images/sdcard.img: 7,23 GiB, 7759462400 bytes, 15155200 sectors
Units: sectors of 1 * 512 = 512 bytes
Sector size (logical/physical): 512 bytes / 512 bytes
I/O size (minimum/optimal): 512 bytes / 512 bytes
Disklabel type: dos
Disk identifier: 0x2df93d98
Device             Boot  Start      End        Sectors   Size  Id  Type
images/sdcard.img1       2048       1026047    1024000   500M  c   W95 FAT32 (LBA)
images/sdcard.img2       1026048    15134719   14108672  6,7G  83  Linux
images/sdcard.img3       15134720   15155199   20480     10M   a2  unknown
\end{shellcodehost}

\subsection{Attach the image to a loop device}

\begin{shellcodehost}
sudo losetup -fP images/sdcard.img
\end{shellcodehost}

Identify the assigned loop device with:

\begin{shellcodehost}
losetup -l
\end{shellcodehost}

Locate the entry corresponding to \repopath{sdcard.img} and note its device name (e.g. \extpath{/dev/loop2}). In each of the commands below, \ccode{N} is to be replaced with the actual number.

\subsection{Format the partitions}

\begin{shellcodehost}
sudo mkfs.vfat -n BOOT /dev/loopNp1
sudo mkfs.ext4 -L rootfs /dev/loopNp2
# loopNp3 is left unformatted, it will receive the raw U-Boot image with dd in the next step
\end{shellcodehost}

\subsection{Mount and copy the files}

All commands below, like the rest of this section, are run from \repopath{sdcard/}; \repopath{boot_partition/}, \repopath{rootfs_overlay/}, \repopath{rootfs.tar} and \repopath{u-boot-with-spl.sfp} are paths relative to it.

\begin{shellcodehost}
cd sdcard/
sudo mkdir -p /tmp/sdcard/boot /tmp/sdcard/rootfs
sudo mount /dev/loopNp1 /tmp/sdcard/boot
sudo mount /dev/loopNp2 /tmp/sdcard/rootfs
\end{shellcodehost}

Copy the boot partition files (kernel, \gterm{devicetree}{DTB}, bitstream, boot script, from Sections~\ref{sec:fpgalix-kernel}, \ref{sec:fpgalix-bitstream} and \ref{sec:fpgalix-boot-script}):

\begin{shellcodehost}
sudo cp -L boot_partition/* /tmp/sdcard/boot/
\end{shellcodehost}

Extract the root filesystem built in Paragraph~\ref{subsec:fpgalix-buildroot-build}:

\begin{shellcodehost}
sudo tar xvf rootfs.tar -C /tmp/sdcard/rootfs/
\end{shellcodehost}

Copy the \gterm{ntp}{NTP}\footnote{\glsentrydesc{ntp}} configuration onto the rootfs:

\begin{shellcodehost}
sudo cp -L rootfs_overlay/etc/chrony.conf /tmp/sdcard/rootfs/etc/
\end{shellcodehost}

Write \gterm{uboot}{U-Boot} (\gterm{spl}{SPL} and \gterm{uboot}{U-Boot} proper, from Section~\ref{sec:fpgalix-bootloader}) raw into the \term{A2} partition:

\begin{shellcodehost}
sudo dd if=u-boot-with-spl.sfp of=/dev/loopNp3 bs=64k seek=0
sync
\end{shellcodehost}

\subsection{Unmount, detach and clean up}

\begin{shellcodehost}
sudo umount /tmp/sdcard/boot /tmp/sdcard/rootfs
sudo losetup -d /dev/loopN
sudo rm -rf /tmp/sdcard
\end{shellcodehost}

\subsection{Write the image to the SD card}

\begin{shellcodehost}
sudo dd if=images/sdcard.img of=/dev/sdX bs=4M status=progress
sync
\end{shellcodehost}

Replace \extpath{/dev/sdX} with the actual \term{SD} card device. \textbf{Verify the target device with particular care} prior to execution, as this command overwrites it in its entirety.

\section{Booting the board}

The \term{MSEL} configuration has already been established in Paragraph~\ref{subsec:fpgalix-msel}.

\begin{enumerate}
    \item Insert the \term{SD} card into the slot on the underside of the board.
    \item Connect a mini-\term{USB} cable between the \term{UART}-to-\term{USB} connector on the board (labelled on the \gterm{pcb}{PCB}) and the \term{PC}.
    \item Connect an \term{Ethernet} cable between the board and the local network.
\end{enumerate}

The board exposes a serial console at 115200 baud, 8N1, no flow control:

\begin{itemize}
    \item \term{Linux}: \cmdpath{picocom}. Identify the device with \lstinline[style=githubcppinline]!dmesg | grep tty! (typically \extpath{/dev/ttyUSB0}), and establish the connection with \cmdpath{picocom -b 115200 /dev/ttyUSB0}.
    \item \term{Windows}: \term{TeraTerm}.\footnote{\term{TeraTerm} is available at \url{https://github.com/TeraTermProject/teraterm}.} Select the \term{COM} port that appears in \term{Device Manager} under \term{Ports (COM \& LPT)} upon connecting the cable.
\end{itemize}

Upon powering on the board, the boot log should appear in the terminal immediately, concluding with a login \gterm{prompt}{prompt}.

\subsection{Logging in and network access}
\label{subsec:fpgalix-login-network}

Log into the board as \sysid{root} over the serial console, with the password set in Paragraph~\ref{subsec:fpgalix-buildroot-base}.

\gterm{ssh}{SSH} access as \sysid{root} does not work, even though the serial console login does: \gterm{openssh}{OpenSSH}'s default \sysid{PermitRootLogin} setting only allows \sysid{root} to log in with a key, not a password, and this project's \sysid{root} account only has a password. Create a new user, therefore, from the serial console:

\begin{shellcodeboard}
mkdir -p /home
adduser <username>
\end{shellcodeboard}

The \cmdpath{mkdir -p /home} is necessary first: \gterm{buildroot}{Buildroot}'s base skeleton does not include a \extpath{/home} directory, and if it is not already present, \gterm{busybox}{BusyBox}'s \cmdpath{adduser} fails silently, leaving the new user without a home directory.

Read the board's \term{IP} address, obtained automatically over \term{Ethernet} by \cmdpath{dhcpcd} (Paragraph~\ref{subsec:fpgalix-buildroot-base}), with either:

\begin{shellcodeboard}
ifconfig
# or
ip addr
\end{shellcodeboard}

From the \term{PC}, verify \gterm{ssh}{SSH} access as the new user:

\begin{shellcodehost}
ssh <username>@<board-ip>
\end{shellcodehost}

